\documentclass[12pt]{article}
\usepackage[T1]{fontenc}
\usepackage{lmodern}
\usepackage{graphicx}
\usepackage{xcolor}
\usepackage{ulem}
\usepackage{setspace}
\setlength{\parindent}{0pt}
\setlength{\parskip}{6pt}
\begin{document}
% setup after \begin{document}: the fast server loads these with the preamble
\newcommand{\reading}[1]{\textbf{#1}\,V}
\renewcommand{\arraystretch}{1.3}

\begin{center}
{\LARGE\bfseries Lab Report 3}\\[4pt]
{\large Measuring Resistance}\\
Henry Yang \quad\textbar\quad 1 October 2026
\end{center}

\noindent\begin{minipage}[t]{0.48\textwidth}
\textbf{Aim.} To measure the resistance of a wire and compare it with the value computed from its resistivity.
\end{minipage}\hfill
\begin{minipage}[t]{0.48\textwidth}
\textbf{Hypothesis.} Resistance is proportional to length, so doubling the length doubles the resistance.
\end{minipage}

\fbox{\parbox{0.9\textwidth}{\textbf{Safety.} Switch off the supply before changing the circuit. \underline{Never} exceed 2\,A.}}

\colorbox{yellow!20}{\textcolor{red}{Important:} record every reading twice, e.g. \reading{1.5} and \reading{1.6}.} Values that were \sout{wrong} corrected are marked \uline{like this}.

\begin{center}
\begin{tabular}{lrr}
Length (m) & Current (A) & Voltage (V) \\
0.5 & 0.40 & 1.1 \\
1.0 & 0.40 & 2.2 \\
\end{tabular}
\end{center}

\begin{quote}
Resistance is futile only when the wire is superconducting. --- \textit{Anonymous}
\end{quote}

\begin{description}
  \item[Voltmeter] connected in parallel across the wire.
  \item[Ammeter] connected in series.
\end{description}

\raggedright
This paragraph is set ragged right with a wide \hspace{1cm} gap and a rule \rule{2cm}{0.4pt} inline, followed by text that goes on long enough to wrap onto the next line of the page.

\centering
\includegraphics[width=0.3\textwidth]{figure.png}

\vspace{1cm}
\begin{spacing}{1.5}
Spaced text at one and a half line spacing, long enough to wrap over two lines so the spacing is visible in the output.
\end{spacing}
\end{document}
